\documentclass[12pt,a4paper]{amsart}

\usepackage{amsmath,amssymb,amsthm}
\usepackage{mathtools}
\usepackage{hyperref}
\usepackage{geometry}
\geometry{margin=2.5cm}

% -------------------------------------------------------
\newtheorem{theorem}{Theorem}[section]
\newtheorem{lemma}[theorem]{Lemma}
\newtheorem{corollary}[theorem]{Corollary}
\newtheorem{proposition}[theorem]{Proposition}
\theoremstyle{definition}
\newtheorem{definition}[theorem]{Definition}
\newtheorem{remark}[theorem]{Remark}

\newcommand{\Z}{\mathbb{Z}}

% -------------------------------------------------------
\begin{document}
% -------------------------------------------------------

\title{No Semiprime is a Giuga Pseudoprime:\\ An Elementary Ordering Argument}

\author{Kurt Ingwer}
\address{Specialist Mathematics Project, Version Build 44}
\email{specialist-maths@localhost}
\date{\today}

\begin{abstract}
A \emph{Giuga pseudoprime} is a composite integer satisfying Giuga's 1950 primality
criterion.  A \emph{semiprime} is a product of exactly two (not necessarily distinct)
primes.  We prove that no semiprime $n = p \cdot q$ with $p < q$ prime is a Giuga
pseudoprime.  The argument is a simple ordering inequality: the weak Giuga condition
applied to the larger prime forces $q \mid (p-1)$, but $p - 1 < q$, so $p - 1 = 0$
and $p = 1$, contradicting primality.  Together with the squarefreeness of Giuga
pseudoprimes, this shows that every Giuga pseudoprime must have at least
\emph{three} distinct prime factors.
\end{abstract}

\maketitle

% -------------------------------------------------------
\section{Introduction}
% -------------------------------------------------------

The Giuga conditions for a composite $n$ with respect to a prime $p \mid n$ are:
\begin{enumerate}
  \item[\textup{(W)}] $p \mid \tfrac{n}{p} - 1$ \quad (weak condition),
  \item[\textup{(S)}] $(p-1) \mid \tfrac{n}{p} - 1$ \quad (strong condition).
\end{enumerate}
A composite $n$ satisfying both conditions for every prime factor is a
\emph{Giuga pseudoprime} \cite{Borwein1996}.

It is already known (and elementary) that every Giuga pseudoprime is squarefree
\cite{Ingwer2026sqfree}.  The present note rules out two-prime composites.

% -------------------------------------------------------
\section{Main Result}
% -------------------------------------------------------

\begin{theorem}
\label{thm:nopq}
There is no Giuga pseudoprime of the form $n = p \cdot q$ with $p < q$ prime.
\end{theorem}

\begin{proof}
Suppose $n = pq$ is a Giuga pseudoprime with primes $p < q$.  The quotients are
$\tfrac{n}{p} = q$ and $\tfrac{n}{q} = p$.

Applying the weak condition~(W) to both prime factors:
\begin{align}
  p &\mid q - 1, \label{eq:Wp} \\
  q &\mid p - 1. \label{eq:Wq}
\end{align}

From~\eqref{eq:Wq}: since $p < q$, we have $0 \le p - 1 < q$.
The only non-negative multiple of~$q$ strictly less than~$q$ is~$0$.
Hence $q \mid (p - 1)$ forces $p - 1 = 0$, so $p = 1$.
But $1$ is not prime — a contradiction.
\end{proof}

\begin{remark}
The argument uses only the weak condition~(W) applied to the \emph{larger} prime~$q$.
Condition~(W) for~$p$ (equation~\eqref{eq:Wp}) is not needed for the contradiction.
\end{remark}

\begin{remark}
The proof is a pure ordering argument: the divisibility $q \mid (p-1)$ combined
with the strict inequality $p - 1 < q$ forces $p - 1 = 0$.  No arithmetic
beyond basic divisibility is required.
\end{remark}

% -------------------------------------------------------
\section{Consequence: At Least Three Prime Factors}
% -------------------------------------------------------

\begin{corollary}
\label{cor:three}
Every Giuga pseudoprime has at least three distinct prime factors.
\end{corollary}

\begin{proof}
A Giuga pseudoprime is squarefree (Theorem~1 in \cite{Ingwer2026sqfree}),
hence a product of distinct primes.  With one prime factor, $n$ would be prime —
not composite.  With two prime factors, Theorem~\ref{thm:nopq} applies.
So at least three are needed.
\end{proof}

\begin{remark}
The lower bound of three is tight in the sense that the conditions do not
immediately rule out three prime factors.  However, a further argument shows
that three factors are also impossible \cite{Ingwer2026a}, so in fact
at least \emph{four} prime factors are required.  The current best lower bound,
due to Bednarek \cite{Bednarek2014}, is $59$ prime factors.
\end{remark}

% -------------------------------------------------------
\section{Symmetry Observation}
% -------------------------------------------------------

If $n = pq$ with $p = q$ (a prime square), then $n$ is not squarefree and already
ruled out.  For $p < q$ the argument above applies.  There is thus a pleasant
symmetry: the two-prime case is eliminated by a combination of
\begin{itemize}
  \item Squarefreeness (eliminates $p = q$), and
  \item The ordering argument (eliminates $p < q$).
\end{itemize}

This two-step exclusion is characteristic of the Giuga theory: structural properties
(squarefreeness) combine with ordering inequalities to yield contradictions.

% -------------------------------------------------------
\begin{thebibliography}{10}

\bibitem{Giuga1950}
G.~Giuga,
\emph{Su una presumibile propriet\`a caratteristica dei numeri primi},
Ist.\ Lombardo Sci.\ Lett.\ Rend.\ Cl.\ Sci.\ Mat.\ Nat.\ \textbf{83}
(1950), 511--528.

\bibitem{Borwein1996}
J.~M.~Borwein, P.~B.~Borwein, R.~Girgensohn, and S.~Parnes,
\emph{Giuga's conjecture on primality},
Amer.\ Math.\ Monthly \textbf{103} (1996), no.\ 1, 40--50.

\bibitem{Bednarek2014}
M.~Bednarek,
\emph{On the size of Giuga pseudoprimes},
Preprint, unpublished (2014).

\bibitem{Ingwer2026sqfree}
K.~Ingwer,
\emph{Every Giuga pseudoprime is squarefree},
Specialist Mathematics Project (2026).

\bibitem{Ingwer2026a}
K.~Ingwer,
\emph{No composite number with exactly three prime factors is a Giuga pseudoprime},
Specialist Mathematics Project (2026).

\end{thebibliography}

\end{document}
